% ============================== CHAPTER 2 ============================
\clearpage
\setlecture{Lecture II}\renewcommand{\tabname}{II}
\chapter{The Chain Rule and the Gradient}\label{ch:chain}

A drone flies through a room where the temperature varies from point to point. The
thermometer on the drone records a single function of time, $T(X(t))$. How fast does that
reading change? It depends on two things: how temperature varies in space (the gradient
of $T$) and how fast and in which direction the drone moves (its velocity). The chain
rule says exactly how to combine them. From it we will get, almost for free, tangent
planes, directional derivatives, the meaning of the gradient, and the conservation of
energy.

\section{The chain rule}

\subsection{Composition along a curve}

Let $f$ be defined on an open set $U$ and let $t\mapsto X(t)$ be a curve with $X(t)\in U$
for every $t$. Then we can form the one-variable function
\[ (f\circ X)(t)=f(X(t)). \]

\begin{example}{a first composition}
Let $f(x,y)=e^x\sin(xy)$ and $X(t)=(t^2,t^3)$. Then $xy=t^5$ and
\[ f(X(t))=e^{t^2}\sin(t^5). \]
Of course we could differentiate this directly with first-year rules. The question is
whether we can get the answer from $\grad f$ and $X'(t)$ alone, without substituting
first. We will come back to this example after the theorem.
\end{example}

\subsection{The dependency picture}

\begin{intuition}
Change travels along paths. When $t$ moves, $x$ moves (at rate $dx/dt$) and $y$ moves
(at rate $dy/dt$). Each of these pushes $f$ (at rates $\partial f/\partial x$ and
$\partial f/\partial y$). Multiply the rates along each path from $t$ to $f$, then add
the paths.
\end{intuition}

\begin{center}
\begin{tikzpicture}[
  v/.style={circle, draw=plum, line width=1pt, fill=plumtint, minimum size=0.95cm, font=\small},
  e/.style={-{Stealth[length=2.4mm]}, line width=1.1pt},
  lb/.style={font=\scriptsize, fill=white, inner sep=1.5pt}]
  \draw[bpgrid, step=0.5] (-0.8,-1.8) grid (8.3,1.8);
  \node[v] (t) at (0,0) {$t$};
  \node[v] (x) at (3.5,1.1) {$x(t)$};
  \node[v] (y) at (3.5,-1.1) {$y(t)$};
  \node[v, draw=coral, fill=coraltint] (f) at (7.2,0) {$f$};
  \draw[e, cyanbp] (t) -- node[lb, above left]{$\dfrac{dx}{dt}$} (x);
  \draw[e, cyanbp] (t) -- node[lb, below left]{$\dfrac{dy}{dt}$} (y);
  \draw[e, coral] (x) -- node[lb, above right]{$\dfrac{\partial f}{\partial x}$} (f);
  \draw[e, coral] (y) -- node[lb, below right]{$\dfrac{\partial f}{\partial y}$} (f);
\end{tikzpicture}

\[
\frac{df}{dt}=\underbrace{\pd fx\frac{dx}{dt}}_{\text{upper path}}+\underbrace{\pd fy\frac{dy}{dt}}_{\text{lower path}}
\]
\end{center}

\begin{theorem}[label=thm:chain]{Chain rule}
Let $f$ be differentiable on an open set $U\subset\R^n$, and let $X:I\to\R^n$ be a
differentiable curve with $X(t)\in U$ for all $t\in I$. Then $t\mapsto f(X(t))$ is
differentiable and
\[
\boxed{\;\frac{d}{dt}f(X(t))=\grad f(X(t))\cdot X'(t)\;}
\]
In components, if $X(t)=(x_1(t),\dots,x_n(t))$,
\[
\frac{d}{dt}f(X(t))=\sum_{i=1}^n\pd f{x_i}(X(t))\,\frac{dx_i}{dt},
\]
and in two variables
\[
\frac{df}{dt}=\pd fx\frac{dx}{dt}+\pd fy\frac{dy}{dt}.
\]
\end{theorem}

The vector form is the clean conceptual statement; the component form is what you use
to compute.

\begin{example}{back to $e^{x}\sin(xy)$ along $(t^2,t^3)$}
\textbf{Step 1: gradient.} $\pd fx=e^x\sin(xy)+e^xy\cos(xy)$ and $\pd fy=e^xx\cos(xy)$.

\textbf{Step 2: velocity.} $X'(t)=(2t,\,3t^2)$.

\textbf{Step 3: chain rule, then substitute $x=t^2$, $y=t^3$, $xy=t^5$.}
\begin{align*}
\frac{d}{dt}f(X(t))&=\big(e^{t^2}\sin t^5+e^{t^2}t^3\cos t^5\big)\,2t+e^{t^2}t^2\cos t^5\cdot 3t^2\\
&=2t\,e^{t^2}\sin t^5+5t^4e^{t^2}\cos t^5 .
\end{align*}
\textbf{Check.} Differentiating $e^{t^2}\sin(t^5)$ directly with the product rule gives
$2te^{t^2}\sin t^5+e^{t^2}\cdot5t^4\cos t^5$. Same answer.
\end{example}

\subsection{Partial derivatives of a composite function}

Very often the inner functions depend on \emph{two} variables:
\[ x=\varphi(t,u),\quad y=\psi(t,u),\qquad g(t,u)=f(\varphi(t,u),\psi(t,u)). \]
To get $\partial g/\partial t$, freeze $u$: then $t\mapsto(\varphi,\psi)$ is a curve and
the chain rule applies. So
\[
\boxed{\;\pd gt=\pd fx\pd xt+\pd fy\pd yt,\qquad \pd gu=\pd fx\pd xu+\pd fy\pd yu.\;}
\]

\begin{center}
\begin{tikzpicture}[
  v/.style={circle, draw=plum, line width=1pt, fill=plumtint, minimum size=0.85cm, font=\small},
  e/.style={-{Stealth[length=2.2mm]}, line width=1pt}]
  \node[v] (t) at (0,0.9) {$t$};
  \node[v] (u) at (0,-0.9) {$u$};
  \node[v] (x) at (3,0.9) {$x$};
  \node[v] (y) at (3,-0.9) {$y$};
  \node[v, draw=coral, fill=coraltint] (f) at (6,0) {$g$};
  \draw[e, cyanbp] (t) -- (x); \draw[e, cyanbp] (t) -- (y);
  \draw[e, cyanbp!40] (u) -- (x); \draw[e, cyanbp!40] (u) -- (y);
  \draw[e, coral] (x) -- (f); \draw[e, coral] (y) -- (f);
  \node[font=\scriptsize, text=navy, align=left, anchor=west] at (6.6,0)
    {for $\partial g/\partial t$:\\ follow the dark paths\\ starting at $t$};
\end{tikzpicture}
\end{center}

\subsection{Proof of the chain rule}

\begin{intuition}
Differentiability says $f$ changes like its linear part $\grad f\cdot K$ when the input
changes by a small vector $K$. Along the curve, the input changes by
$K=X(t+h)-X(t)\approx hX'(t)$. So $f$ changes by about $\grad f\cdot hX'(t)$. Divide by
$h$ and the error term disappears in the limit.
\end{intuition}

\begin{proof}
\textbf{Step 1: expand the increment of $f$.} Fix $t$ and set $K=K(h)=X(t+h)-X(t)$.
Since $f$ is differentiable at $X(t)$,
\[ f(X(t)+K)-f(X(t))=\grad f(X(t))\cdot K+\norm K\,g(K),\qquad g(K)\to0\text{ as }K\to O, \]
where we set $g(O)=0$ (this handles the values of $h$ for which $K=O$).

\textbf{Step 2: divide by $h$.} Since $X(t)+K=X(t+h)$,
\begin{align*}
\frac{f(X(t+h))-f(X(t))}{h}&=\grad f(X(t))\cdot\frac{X(t+h)-X(t)}{h}\\
&\quad\pm\ \norm{\frac{X(t+h)-X(t)}{h}}\,g(K).
\end{align*}
The sign is $+$ if $h>0$ and $-$ if $h<0$, because $\norm K/h=\pm\norm{K/h}$.

\textbf{Step 3: let $h\to0$.} The curve is differentiable, so
$\frac{X(t+h)-X(t)}{h}\to X'(t)$; in particular its norm stays bounded. The curve is
also continuous, so $K\to O$ and hence $g(K)\to0$. The last term is (bounded)$\times$(tends
to $0$), so it tends to $0$. Therefore
\[ \frac{d}{dt}f(X(t))=\grad f(X(t))\cdot X'(t). \]
\end{proof}

\subsection{Worked examples}

\begin{example}{polar coordinates}
Let $f(x,y)=x^2+2xy$ with $x=r\cos\theta$, $y=r\sin\theta$, and $g(r,\theta)=f(r\cos\theta,r\sin\theta)$.
Find $\partial g/\partial\theta$.

\textbf{Step 1.} $\pd fx=2x+2y$, $\pd fy=2x$, $\pd x\theta=-r\sin\theta$, $\pd y\theta=r\cos\theta$.

\textbf{Step 2.}
\[ \pd g\theta=(2x+2y)(-r\sin\theta)+2x(r\cos\theta). \]
\textbf{Step 3: substitute and simplify.}
\begin{align*}
\pd g\theta&=-2r^2\cos\theta\sin\theta-2r^2\sin^2\theta+2r^2\cos^2\theta
 = r^2\big(2\cos2\theta-\sin2\theta\big).
\end{align*}
\textbf{Check.} $g=r^2\cos^2\theta+r^2\sin2\theta$, so directly
$\pd g\theta=-r^2\sin2\theta+2r^2\cos2\theta$. Same.
\end{example}

\begin{example}{a three-variable composition}
Let $w=f(x,y,z)=e^{xy}\cos z$ with $x=tu$, $y=\sin(tu)$, $z=u^2$. Find $\partial w/\partial u$.

Three paths from $u$ to $w$:
\[ \pd wu=\pd fx\pd xu+\pd fy\pd yu+\pd fz\pd zu . \]
With $\pd fx=ye^{xy}\cos z$, $\pd fy=xe^{xy}\cos z$, $\pd fz=-e^{xy}\sin z$ and
$\pd xu=t$, $\pd yu=t\cos(tu)$, $\pd zu=2u$:
\[
\pd wu=ye^{xy}(\cos z)\,t+xe^{xy}(\cos z)\cos(tu)\,t-e^{xy}(\sin z)\,2u .
\]
Substituting $x=tu$, $y=\sin(tu)$, $z=u^2$:
\begin{align*}
\pd wu&=t\sin(tu)\,e^{tu\sin(tu)}\cos(u^2)+t^2u\,e^{tu\sin(tu)}\cos(u^2)\cos(tu)\\
&\quad-2u\,e^{tu\sin(tu)}\sin(u^2).
\end{align*}
\end{example}

\begin{example}{be careful with variable names}
Let $u=f(x^2-y,\,xy)$ where $f$ is an unknown differentiable function. Find
$\partial u/\partial x$ and $\partial u/\partial y$.

The letters $x,y$ are already the outer variables, so we must not call the slots of $f$
by those names. Introduce
\[ s=x^2-y,\qquad t=xy . \]
Then $\pd sx=2x$, $\pd tx=y$, $\pd sy=-1$, $\pd ty=x$, and
\begin{align*}
\pd ux&=\pd fs\,2x+\pd ft\,y=2x\,D_1f(s,t)+y\,D_2f(s,t),\\
\pd uy&=-\pd fs+x\pd ft=-D_1f(s,t)+x\,D_2f(s,t).
\end{align*}

\textbf{Going backwards.} Writing $u=g(x,y)$, we can solve these two equations for the
partials of $f$. Multiply the second by $2x$ and add it to the first:
\[
\pd gx+2x\pd gy=(y+2x^2)\pd ft\quad\Longrightarrow\quad
\pd ft=\frac{1}{y+2x^2}\left(\pd gx+2x\pd gy\right).
\]
\end{example}

\begin{trap}
When the outer variables already use the letters $x,y$, writing $\partial f/\partial x$
for ``the derivative of $f$ in its first slot'' is ambiguous. Use $D_1f$, $D_2f$, or give
the slots new names ($s,t$) as above.
\end{trap}

\begin{example}{a scalar inner function}
Let $f(x,y,z)=g(x^2-3zy+xz)$ with $g$ a differentiable function of one variable. With
$h(x,y,z)=x^2-3zy+xz$, the one-variable chain rule gives
\[
\pd fx=g'(h(x,y,z))\pd hx=g'(x^2-3zy+xz)\,(2x+z).
\]
Similarly $\pd fy=g'(h)\,(-3z)$ and $\pd fz=g'(h)\,(x-3y)$. In vector form:
$\grad f=g'(h)\,\grad h$.
\end{example}

\begin{example}{$D_1f$, $D_2f$ notation}
Let $g(t,x,y)=f(t^2x,\,ty)$. Then
\[
\pd gt=D_1f(t^2x,ty)\cdot 2tx+D_2f(t^2x,ty)\cdot y .
\]
Here $x$ is an outer variable, so ``$\partial f/\partial x$'' would be ambiguous.
\end{example}

\begin{example}{temperature along a flight path}
The temperature field is $T(x,y)=20+2x-y^2$ and a drone follows $X(t)=(t,t^2)$. Is the
drone warming up or cooling down at $t=1$?

$\grad T(x,y)=(2,-2y)$ and $X'(t)=(1,2t)$. At $t=1$ the drone is at $X(1)=(1,1)$, so
\[
\frac{d}{dt}T(X(t))\Big|_{t=1}=\grad T(1,1)\cdot X'(1)=(2,-2)\cdot(1,2)=2-4=-2 .
\]
The reading drops at $2$ degrees per unit time: the drone is flying into colder air.
\end{example}

\begin{reallife}{backpropagation}
Training a neural network means computing the derivative of a loss with respect to
millions of weights, through many layers of composed functions. Backpropagation is
nothing but the chain rule applied along the dependency graph, from the output back to
the inputs, exactly like the path diagrams above.
\end{reallife}

\subsection{A first geometric consequence}

\begin{example}{a path of constant temperature}
Let $f:\R^3\to\R$ be a temperature and $t\mapsto B(t)$ the path of a particle that
always stays at the same temperature: $f(B(t))=c$. Differentiating the constant,
\[ \grad f(B(t))\cdot B'(t)=0 . \]
So the gradient is perpendicular to the velocity of every constant-temperature path.
\end{example}

\section{Tangent planes}

\subsection{Level surfaces and the gradient}

Let $f$ be differentiable and $c$ a constant. Consider the level surface
\[ f(X)=c,\qquad\text{with }\grad f(X)\neq O . \]
If a differentiable curve $X(t)$ lies on the surface, then $f(X(t))=c$ for all $t$, and
by the chain rule
\[ \grad f(X(t))\cdot X'(t)=0 . \]
Every tangent direction $X'(t)$ of the surface is perpendicular to the gradient.

\begin{keyidea}
\textbf{The gradient is perpendicular to level curves and level surfaces.} It crosses them
at right angles, like the steepest path down a hill crosses the contour lines.
\end{keyidea}

\begin{center}
\begin{tikzpicture}
\begin{axis}[bp, width=7.4cm, height=6.2cm, xmin=-0.2, xmax=3.6, ymin=-0.2, ymax=3.0,
   xtick=\empty, ytick=\empty, xlabel=$x$, ylabel=$y$]
  \foreach \c in {1.5,3,4.5}
    {\edef\tmp{\noexpand\addplot[plum!70, domain=0.25:3.5, samples=60] {\c/x};}\tmp}
  \addplot[plum, line width=1.6pt, domain=0.6:3.5, samples=60] {3/x};
  \addplot[cyanbp, dashed, line width=1pt, domain=0.6:3.0] {-1.3333*(x-1.5)+2};
  \draw[vec] (axis cs:1.5,2) -- (axis cs:2.3,2.6) node[right, font=\small]{$\nabla f(P)$};
  \node[dot] at (axis cs:1.5,2) {};
  \node[font=\small, below left, text=plumdark] at (axis cs:1.5,2) {$P$};
  \node[font=\scriptsize, text=cyanbp, anchor=west] at (axis cs:2.4,0.3) {tangent};
  \node[font=\scriptsize, text=plum, fill=white, inner sep=1pt] at (axis cs:3.2,1.25) {$f=c$};
\end{axis}
\end{tikzpicture}

{\small\color{navy} Level curves of $f(x,y)=xy$. At $P$ the gradient $(y,x)$ is
perpendicular to the tangent of the level curve through $P$.}
\end{center}

\subsection{Definition and examples}

\begin{definition}{Tangent plane}
Let $P$ be a point on the level surface $f(X)=c$, with $\grad f(P)\neq O$. The
\emph{tangent plane} to the surface at $P$ is the plane through $P$ perpendicular to
$N=\grad f(P)$:
\[
 \boxed{\,X\cdot N=P\cdot N\,}\qquad\text{equivalently}\qquad (X-P)\cdot\grad f(P)=0 .
\]
In the plane ($n=2$) the same equation gives the \emph{tangent line} to a level curve.
\end{definition}

\begin{recipe}
To find the tangent plane to a surface at $P$:
\begin{enumerate}
  \item Write the surface as a level set $f(x,y,z)=c$ (move everything to one side).
  \item Compute $N=\grad f(P)$.
  \item Write $N\cdot(X-P)=0$ and simplify.
\end{enumerate}
\end{recipe}

\begin{example}{tangent plane to a sphere}
Find the tangent plane to $x^2+y^2+z^2=3$ at $P=(1,1,1)$.

$f=x^2+y^2+z^2$, $\grad f=(2x,2y,2z)$, $N=\grad f(P)=(2,2,2)$. Then $X\cdot N=P\cdot N$
reads $2x+2y+2z=6$, that is
\[ x+y+z=3 . \]
As expected, the normal $(1,1,1)$ points along the radius: tangent planes to a sphere are
perpendicular to the radius.
\end{example}

\begin{example}{tangent line to a level curve}
Find the tangent line to $x^2y+y^3=10$ at $P=(1,2)$.

Check that $P$ is on the curve: $1\cdot2+8=10$. With $f=x^2y+y^3$,
$\grad f=(2xy,\;x^2+3y^2)$, so $N=\grad f(1,2)=(4,13)$. The tangent line is
$4x+13y=4\cdot1+13\cdot2$, that is
\[ 4x+13y=30 . \]
\end{example}

\begin{center}
\begin{tikzpicture}
\begin{axis}[bp, width=7.6cm, height=6.6cm, xmin=-3.6, xmax=3.6, ymin=-0.3, ymax=3.3,
   xtick={-3,-2,-1,1,2,3}, ytick={1,2,3}, xlabel=$x$, ylabel=$y$]
  \addplot[plum, line width=1.5pt, domain=0.6:2.154, samples=80] ({sqrt(max(0,(10-x^3)/x))},{x});
  \addplot[plum, line width=1.5pt, domain=0.6:2.154, samples=80] ({-sqrt(max(0,(10-x^3)/x))},{x});
  \addplot[cyanbp, dashed, line width=1pt, domain=-1.6:3.5] {(30-4*x)/13};
  \draw[vec] (axis cs:1,2) -- (axis cs:1.28,2.91) node[right, font=\small]{$(4,13)$};
  \node[dot] at (axis cs:1,2) {};
  \node[font=\small, below left, text=plumdark] at (axis cs:1,2) {$P$};
  \node[font=\scriptsize, text=plum, fill=white, inner sep=1pt] at (axis cs:-2.5,1.4) {$x^2y+y^3=10$};
  \node[font=\scriptsize, text=cyanbp, fill=white, inner sep=1pt] at (axis cs:2.7,1.65) {$4x+13y=30$};
\end{axis}
\end{tikzpicture}
\end{center}

\subsection{Surfaces given as graphs}

A graph $z=g(x,y)$ can always be written as a level surface:
\[ f(x,y,z)=g(x,y)-z=0,\qquad \grad f=\big(g_x,\;g_y,\;-1\big). \]
The tangent plane at $P=(a,b,g(a,b))$ is $g_x(a,b)(x-a)+g_y(a,b)(y-b)-(z-g(a,b))=0$, that is
\[
\boxed{\,z=g(a,b)+g_x(a,b)(x-a)+g_y(a,b)(y-b)\,}
\]
This is exactly the linear part from the definition of differentiability (Chapter 1).

\begin{example}{tangent plane to a paraboloid}
For $z=x^2+y^2$: $f=x^2+y^2-z$, $\grad f=(2x,2y,-1)$. At $P=(1,2,5)$,
$\grad f(P)=(2,4,-1)$, and $X\cdot N=P\cdot N$ gives
\[ 2x+4y-z=2+8-5=5 . \]
At $P=(1,1,2)$ the same method gives $2x+2y-z=2$, i.e. $z=2x+2y-2$ (picture below).
\end{example}

\begin{center}
\begin{tikzpicture}
\begin{axis}[bp3d, width=10cm, height=8.4cm, xlabel=$x$, ylabel=$y$, zlabel=$z$,
   xmin=-0.5, xmax=2.4, ymin=-0.5, ymax=2.4, zmin=-1, zmax=8.5, view={-48}{20}]
  \addplot3[plumsurf, domain=0:2, y domain=0:2, samples=16, opacity=0.85, fill opacity=0.85] {x^2+y^2};
  \addplot3[surf, shader=flat, draw=cyanbp, fill=cyanbp, fill opacity=0.3, draw opacity=0.8,
     domain=0.35:1.65, y domain=0.35:1.65, samples=5] {2*x+2*y-2};
  \addplot3[-{Stealth[length=2.5mm]}, coral, line width=1.3pt] coordinates {(1,1,2) (1.6,1.6,1.7)};
  \addplot3[only marks, mark=*, mark size=2.6pt, coral] coordinates {(1,1,2)};
  \node[font=\scriptsize, text=coraldark, anchor=west] at (axis cs:1.6,1.6,1.4) {$N=(2,2,-1)$};
  \node[font=\scriptsize, text=cyanbp, anchor=east] at (axis cs:0.35,0.35,-0.6) {$z=2x+2y-2$};
\end{axis}
\end{tikzpicture}
\end{center}

\begin{extra}
\textbf{Linear approximation in action.} Estimate $\sqrt{3.02^2+3.97^2}$ without a calculator.
Take $g(x,y)=\sqrt{x^2+y^2}$ near $(3,4)$: $g(3,4)=5$, $\grad g=\left(\frac xg,\frac yg\right)=\left(\frac35,\frac45\right)$,
and $H=(0.02,-0.03)$. Then
\[
g(3.02,3.97)\approx5+\tfrac35(0.02)+\tfrac45(-0.03)=5+0.012-0.024=4.988 .
\]
The true value is $4.98811\ldots$, correct to four digits.
\end{extra}

\section{Directional derivatives}

\begin{intuition}
Partial derivatives give slopes due east and due north. But you can walk in any direction.
The slope in the direction of a unit vector $A$ is obtained by walking along the straight
line $P+tA$ and differentiating at $t=0$. The chain rule computes it immediately.
\end{intuition}

Let $P$ be a point and $A$ a \emph{unit} vector, $\norm A=1$. The line through $P$ in
direction $A$ is $X(t)=P+tA$, with $X'(t)=A$. By the chain rule,
$\frac{d}{dt}f(P+tA)=\grad f(P+tA)\cdot A$.

\begin{definition}{Directional derivative}
The \emph{directional derivative} of $f$ at $P$ in the direction of the unit vector $A$ is
\[
\boxed{\,D_Af(P)=\frac{d}{dt}f(P+tA)\Big|_{t=0}=\grad f(P)\cdot A\,}
\]
It is the rate of change of $f$ at $P$ along the straight line through $P$ in direction $A$.
\end{definition}

\mnote{For $A=E_1=(1,0,\dots,0)$, $D_Af=D_1f$. Partial derivatives are directional
derivatives along the axes.}

\begin{trap}
\textbf{Normalise first.} If the direction is given by a non-unit vector $B$, use
$A=B/\norm B$. Forgetting this multiplies the answer by $\norm B$, and it is the most
common error on this topic.
\end{trap}

\begin{example}{a directional derivative}
Let $f(x,y)=x^2+y^3$, $P=(-1,3)$ and direction $B=(1,2)$.

\textbf{Step 1: normalise.} $\norm B=\sqrt5$, so $A=\frac{1}{\sqrt5}(1,2)$.

\textbf{Step 2: gradient.} $\grad f=(2x,3y^2)$, so $\grad f(P)=(-2,27)$.

\textbf{Step 3: dot product.}
\[ D_Af(P)=(-2,27)\cdot\tfrac{1}{\sqrt5}(1,2)=\frac{-2+54}{\sqrt5}=\frac{52}{\sqrt5}. \]
\end{example}

\subsection{Geometry: the gradient is the direction of steepest ascent}

For a unit vector $A$, let $\theta$ be the angle between $A$ and $\grad f(P)$. Then
\[ D_Af(P)=\grad f(P)\cdot A=\norm{\grad f(P)}\,\norm A\cos\theta=\norm{\grad f(P)}\cos\theta . \]

\begin{center}
\begin{tikzpicture}[scale=1.0]
  \draw[bpgrid, step=0.5] (-2.6,-1.6) grid (3.6,2.6);
  \draw[vec] (0,0) -- (3,1.5) node[right, font=\small]{$\nabla f(P)$};
  \draw[vecc] (0,0) -- (0.6,1.4) node[above, font=\small]{$A$};
  \draw[navy, line width=0.6pt] (26.57:0.7) arc (26.57:66.8:0.7);
  \node[font=\small, text=navy] at (47:0.95) {$\theta$};
  \draw[navy, dashed] (0.6,1.4) -- (1.16,0.58);
  \draw[line width=2.5pt, cyanbp!45] (0,0) -- (1.16,0.58);
  \node[font=\scriptsize, text=cyanbp, rotate=26.57] at (0.62,0.08) {$D_Af(P)$ (scaled)};
  \node[dot] at (0,0) {};
  \node[font=\small, below left] at (0,0) {$P$};
  \draw[plum!60, dashed, line width=1pt] (-2.2,1.1) -- (2.2,-1.1) node[right, font=\scriptsize, text=plum]{level curve};
\end{tikzpicture}
\end{center}

\begin{itemize}
  \item $\theta=0$ ($A$ along $\grad f$): $D_Af=\norm{\grad f}$, the \textbf{maximal increase}.
  \item $\theta=\pi/2$ ($A$ along the level curve): $D_Af=0$, no first-order change.
  \item $\theta=\pi$ ($A$ against $\grad f$): $D_Af=-\norm{\grad f}$, the \textbf{maximal decrease}.
\end{itemize}

\begin{theorem}{Meaning of the gradient}
Assume $\grad f(P)\neq O$. For every unit vector $A$,
\[ -\norm{\grad f(P)}\le D_Af(P)\le\norm{\grad f(P)} . \]
The maximum is attained exactly for $A=\dfrac{\grad f(P)}{\norm{\grad f(P)}}$ and the
minimum for the opposite vector. So the gradient points in the direction of maximal
increase of $f$, and its length is that maximal rate.
\end{theorem}

\begin{proof}
From $D_Af(P)=\norm{\grad f(P)}\cos\theta$ and $-1\le\cos\theta\le1$ we get the two
inequalities. Equality on the right holds if and only if $\cos\theta=1$, i.e. $\theta=0$,
i.e. $A$ is the unit vector in the direction of $\grad f(P)$. In that case indeed
\[ D_Af(P)=\grad f(P)\cdot\frac{\grad f(P)}{\norm{\grad f(P)}}=\frac{\norm{\grad f(P)}^2}{\norm{\grad f(P)}}=\norm{\grad f(P)} . \]
Equality on the left holds iff $\theta=\pi$, i.e. $A=-\grad f(P)/\norm{\grad f(P)}$.
\end{proof}

\begin{example}{which way up the hill?}
A hill has height $h(x,y)=100-x^2-2y^2$. A hiker stands at $P=(2,1)$.

\textbf{Steepest ascent.} $\grad h=(-2x,-4y)$, so $\grad h(P)=(-4,-4)$, with
$\norm{\grad h(P)}=4\sqrt2$. The best direction is
\[ A_{\max}=\frac{1}{\sqrt2}(-1,-1),\qquad D_{A_{\max}}h(P)=4\sqrt2\approx5.66 . \]
\textbf{Walking east.} With $E=(1,0)$: $D_Eh(P)=(-4,-4)\cdot(1,0)=-4$. Going east is
downhill.

\textbf{Walking along the contour.} Directions perpendicular to $(-1,-1)$, e.g.
$\frac{1}{\sqrt2}(1,-1)$, give $D_Ah(P)=0$: the hiker stays at the same height to first order.
\end{example}

\begin{center}
\begin{tikzpicture}
\begin{axis}[bp, width=7.6cm, height=6.4cm, xmin=-3.2, xmax=3.2, ymin=-2.3, ymax=2.3,
   xtick={-2,2}, ytick={-1,1}, xlabel=$x$, ylabel=$y$, axis equal image]
  \foreach \c in {1,2,3,4,5,6,7,8} {
    \edef\tmp{\noexpand\addplot[plum!\the\numexpr 30+8*\c\relax!coral, domain=0:360, samples=80]
      ({sqrt(\c)*cos(x)},{sqrt(\c/2)*sin(x)});}\tmp }
  \node[dot] at (axis cs:2,1) {};
  \draw[vec] (axis cs:2,1) -- (axis cs:1.2,0.2);
  \node[font=\scriptsize, text=coraldark, anchor=east, fill=white, inner sep=1pt] at (axis cs:1.15,0.15) {steepest};
  \draw[vecc] (axis cs:2,1) -- (axis cs:2.9,1);
  \node[font=\scriptsize, text=cyanbp, anchor=south] at (axis cs:2.75,1.05) {east};
\end{axis}
\end{tikzpicture}

{\small\color{navy} Contours of $100-x^2-2y^2$ (ellipses). The steepest direction at $P$ is
perpendicular to the contour through $P$.}
\end{center}

\begin{reallife}{gradient descent}
To minimise a cost function $C(w_1,\dots,w_n)$, move a small step \emph{against} the
gradient: $w\leftarrow w-\eta\,\grad C(w)$. By the theorem above, $-\grad C$ is the
direction in which $C$ decreases fastest. This one line is how almost every machine
learning model is trained. A heat-seeking robot does the opposite: it climbs $+\grad T$.
\end{reallife}

\section{The conservation law}

\begin{definition}{Vector field}
A \emph{vector field} on an open set $U$ associates to each point $X\in U$ a vector
$F(X)$ of the same dimension. If $f$ is differentiable, $X\mapsto\grad f(X)$ is a vector
field (the \emph{gradient field} of $f$).
\end{definition}

\begin{definition}{Conservative field}
A vector field $F$ is \emph{conservative} if there is a differentiable function $\Phi$,
called a \emph{potential}, with
\[ F=-\grad\Phi . \]
\end{definition}

The minus sign is a physics convention: forces push from high potential energy to low, that
is, downhill on $\Phi$.

\subsection{From Newton's law to energy}

\begin{intuition}
Kinetic energy grows when the force pushes you along your motion, and the force pushes
you downhill on $\Phi$. Whatever kinetic energy you gain, you lose in potential energy.
The chain rule is the bookkeeping that makes this exact.
\end{intuition}

\begin{theorem}{Conservation of mechanical energy}
A particle of mass $m$ moves along $X(t)$ under a conservative force $F=-\grad\Phi$,
following Newton's law $mX''(t)=F(X(t))$. Then
\[
\underbrace{\tfrac12m\norm{X'(t)}^2}_{\text{kinetic energy}}+\underbrace{\Phi(X(t))}_{\text{potential energy}}=\text{constant}.
\]
\end{theorem}

\begin{proof}
\textbf{Step 1.} Newton's law with $F=-\grad\Phi$ reads $mX''+\grad\Phi(X)=O$.

\textbf{Step 2.} Take the dot product with the velocity $X'$:
\[ mX'\cdot X''+\grad\Phi(X)\cdot X'=0 . \]
\textbf{Step 3: recognise two derivatives.} Since $\norm{X'}^2=X'\cdot X'$, the product
rule for dot products gives
\[ \frac{d}{dt}\left(\tfrac12m\norm{X'(t)}^2\right)=\tfrac12m\,(2X'\cdot X'')=mX'\cdot X'' . \]
By the chain rule, $\frac{d}{dt}\Phi(X(t))=\grad\Phi(X(t))\cdot X'(t)$.

\textbf{Step 4.} So Step 2 says
$\frac{d}{dt}\left(\tfrac12m\norm{X'}^2+\Phi(X)\right)=0$, and a function with zero
derivative on an interval is constant.
\end{proof}

\subsection{The spring}

\begin{example}{mass-spring oscillator}
For a spring, $F(x)=-kx=-\frac{d}{dx}\left(\tfrac12kx^2\right)$, so $\Phi(x)=\tfrac12kx^2$.
Newton's law is $m\ddot x+kx=0$. Multiply by $\dot x$:
\[ m\dot x\ddot x+kx\dot x=\frac{d}{dt}\left(\tfrac12m\dot x^2+\tfrac12kx^2\right)=0, \]
hence $\tfrac12m\dot x^2+\tfrac12kx^2=E$ is constant. At the turning points $\pm x_{\max}$
the speed is zero, so $E=\tfrac12kx_{\max}^2$, and the maximum speed (at $x=0$) is
$\dot x_{\max}=x_{\max}\sqrt{k/m}$.
\end{example}

\begin{center}
\begin{tikzpicture}
\begin{axis}[bp, width=8cm, height=5.6cm, xmin=-2.6, xmax=2.6, ymin=-0.3, ymax=3.4,
   xtick={-2,2}, xticklabels={$-x_{\max}$,$x_{\max}$}, ytick=\empty,
   xlabel=$x$, ylabel={energy}]
  \addplot[name path=phi, plum, line width=1.5pt, domain=-2.45:2.45, samples=60] {0.5*x^2};
  \addplot[name path=E, coral, line width=1.2pt, domain=-2:2] {2};
  \addplot[cyanbp, fill opacity=0.25, fill=cyanbp, draw=none] fill between[of=E and phi, soft clip={domain=-2:2}];
  \draw[navy, dashed, line width=0.7pt] (axis cs:-2,0) -- (axis cs:-2,2);
  \draw[navy, dashed, line width=0.7pt] (axis cs:2,0) -- (axis cs:2,2);
  \node[font=\scriptsize, text=coraldark, anchor=south] at (axis cs:0,2.05) {total energy $E$};
  \node[font=\scriptsize, text=cyanbp] at (axis cs:0,1.1) {kinetic};
  \node[font=\scriptsize, text=plum, anchor=west] at (axis cs:1.6,2.9) {$\Phi=\tfrac12kx^2$};
\end{axis}
\end{tikzpicture}

{\small\color{navy} The gap between the line $E$ and the curve $\Phi$ is the kinetic
energy. It vanishes at the turning points.}
\end{center}

\subsection{Inverse-square forces}

Consider a force directed along the position vector, with magnitude inversely
proportional to the square of the distance $r=\norm X$:
\[ F(X)=\frac{C}{\norm X^2}\,\frac{X}{\norm X}=C\,\frac{X}{r^3}. \]

\textbf{Finding a potential.} First compute $\grad(1/r)$. Since
$r=(x^2+y^2+z^2)^{1/2}$,
\[
\pd{}{x}\left(\frac1r\right)=-\frac12(x^2+y^2+z^2)^{-3/2}\cdot2x=-\frac{x}{r^3},
\qquad\text{so}\qquad \grad\left(\frac1r\right)=-\frac{X}{r^3}.
\]
Therefore, with $\Phi(X)=\dfrac{C}{r}$,
\[ -\grad\Phi=-C\grad\left(\frac1r\right)=C\frac{X}{r^3}=F . \]
So $F$ is conservative with potential $\Phi=C/r$.

\begin{trap}
\textbf{Check the sign by computing $-\grad\Phi$.} For gravity the force points
\emph{towards} the origin, so $C=-GMm<0$ and $\Phi=-\dfrac{GMm}{r}$, the familiar formula.
This is the ``$\Phi=-C/r$'' you may see written with $C>0$ standing for the strength of an
attractive force, $F=-C X/r^3$. Both statements agree once the sign of the force is fixed.
\end{trap}

\begin{reallife}{escape velocity}
Energy conservation with $\Phi=-GMm/r$ gives the escape velocity from a planet of radius
$R$: starting with speed $v$ and arriving at $r=\infty$ with speed $0$ requires
$\tfrac12mv^2-\frac{GMm}{R}=0$, so $v=\sqrt{2GM/R}$. For the Earth this is about
$11.2$ km/s, computed with no differential equation solved at all.
\end{reallife}

\section{Summary}

\begin{keyidea}
\begin{itemize}
  \item Chain rule: $\dfrac{d}{dt}f(X(t))=\grad f(X(t))\cdot X'(t)$; add one term per path in the dependency graph.
  \item The gradient is perpendicular to level curves and surfaces.
  \item Tangent plane to $f=c$ at $P$: $(X-P)\cdot\grad f(P)=0$; for a graph, $z=g+g_x(x-a)+g_y(y-b)$.
  \item Directional derivative: $D_Af(P)=\grad f(P)\cdot A$ with $\norm A=1$.
  \item $\grad f$ points to steepest ascent; its length is the steepest slope.
  \item Conservative forces $F=-\grad\Phi$ conserve kinetic plus potential energy.
\end{itemize}
\end{keyidea}

\section{Exam-style questions}

\begin{exercise}[1]
Let $f(x,y)=x^2y-y^3$ and $X(t)=(\cos t,\sin t)$. Compute $\frac{d}{dt}f(X(t))$ using the
chain rule, then check by substituting first.
\end{exercise}

\begin{exercise}[1]
Find the tangent plane to the surface $z=xe^{y}$ at the point $(2,0,2)$.
\end{exercise}

\begin{exercise}[2]
Let $f(x,y,z)=xy+yz+zx$. Find the directional derivative of $f$ at $P=(1,-1,2)$ in the
direction of $B=(2,1,-2)$. In which direction does $f$ increase fastest at $P$, and what
is the maximal rate?
\end{exercise}

\begin{exercise}[2]
Let $z=f(x,y)$ with $x=r\cos\theta$, $y=r\sin\theta$. Show that
\[
\left(\pd fx\right)^2+\left(\pd fy\right)^2=\left(\pd zr\right)^2+\frac{1}{r^2}\left(\pd z\theta\right)^2 .
\]
\end{exercise}

\begin{exercise}[2]
Find all points on the surface $x^2+2y^2+3z^2=12$ where the tangent plane is parallel to
the plane $x+4y+6z=0$.
\end{exercise}

\begin{exercise}[2]
Let $u=f(x-y,\,y-z,\,z-x)$ where $f$ is differentiable. Show that
$\pd ux+\pd uy+\pd uz=0$.
\end{exercise}

\begin{exercise}[3]
The temperature in space is $T(x,y,z)=e^{-x^2-2y^2-3z^2}$. A bee at $P=(1,1,1)$ wants to
cool down as fast as possible. In which direction should it fly? If it flies with velocity
$X'=(1,0,0)$ instead, at what rate does the temperature it feels change?
\end{exercise}

\begin{exercise}[3]
A particle of mass $m=2$ moves in the plane under the force $F(x,y)=(-2x,\,-8y)$.
\begin{enumerate}[label=(\alph*)]
  \item Find a potential $\Phi$ with $F=-\grad\Phi$.
  \item The particle starts at $(1,0)$ with velocity $(0,3)$. Use conservation of energy
  to find its speed whenever it passes through the point $(0,1)$.
  \item Explain why the particle can never reach the point $(0,2)$.
\end{enumerate}
\end{exercise}
